% GENERATED, edit the French unit or the translation table
% GENERATED from content/data/enrichment.json by mxbuild.enrichment; do not edit.
\mxunit{enrichment-bisection-notes}{Bracketing a root by bisection}
\mxsubtitle{Upper high school · analysis and algorithms \textperiodcentered{} Lesson and method \textperiodcentered{} 12}

\begin{prerequisites}
Continuity, signs and intervals.
\end{prerequisites}

\begin{keypoint}[Ideas and method]
If f is continuous on [a,b] and f(a), f(b) have opposite signs, the intermediate value theorem guarantees at least one root in [a,b]. It does not guarantee uniqueness.

Compute $m=\frac{a+b}{2}$. If $f(m)=0$, a root has been found. Otherwise retain the half whose endpoint values have opposite signs. Continuity maintains the guarantee at every step.

After n halvings, the interval width is $\frac{b-a}{2^{n}}$. Its midpoint approximates a root with error at most half that width. Width is therefore different from midpoint error.
\end{keypoint}

\begin{example}
For $f(x)=x^{2}-2$ on $[1 ; 2]$, f is continuous and strictly increasing, so its root is unique. Midpoints $1.5$ and then $1.25$ give intervals $[1 ; 1.5]$ and $[1.25 ; 1.5]$: the left value remains negative and the right value positive.
\end{example}

\begin{activity}[Explore]
On a number line, draw the successive intervals around $\sqrt{2}$ and label each width.
\end{activity}

